\section{Exploración segura y operación prolongada}
\subsection{Optimización con restricciones y límites suaves}
\begin{figure}[H]
\centering
\includegraphics[width=0.9\textwidth]{slides-images/slide-06.jpg}
\caption{La Slide 6 trata las restricciones de seguridad como una reward penalty que se usa en el soft policy gradient.}
\end{figure}

\begin{knowledgebox}{Estrategias de exploración segura}
\begin{itemize}
\item Convertir la safety constraint en un penalty term $c(s_t,a_t)$;\\
\item Controlar la intensidad de la penalty con el Lagrangian multiplier $\lambda$;\\
\item Cuando la uncertainty es alta (por ejemplo, ante un new goal), reducir el exploration rate para evitar acciones high-risk.
\end{itemize}
\end{knowledgebox}

\subsection{Monitoreo con un horizon prolongado}
Chelsea enfatiza que la Autonomy no es un único entrenamiento, sino una property de múltiples deployments; por eso se requiere un monitoreo a nivel de sistema:
\begin{itemize}
\item Registrar las trayectorias de success/failure de cada goal;
\item Hacer track de $\Delta reward$ y $\Delta constraint$ para detectar drifting;
\item Ante fallas repetidas, revertir dinámicamente el modelo de reward o reducir el goal space.
\end{itemize}

\begin{warningbox}{Conservar una vía de reversión en el despliegue}
El pipeline de despliegue debe conservar los mecanismos de «revertir a un checkpoint» y «reentrenar la estimación de reward», para poder resolver con rapidez problemas como una hallucination reward o una reset policy que deja de funcionar.
\end{warningbox}

\subsection{Resumen de la sección}
La exploración segura protege el comportamiento mediante una arquitectura de penalty/constraint; la operación prolongada depende de estadísticas de monitoreo y de planes de reversión, y en conjunto conducen a un sistema de Autonomy confiable.
\newpage
